\documentclass[11pt,letterpaper]{article}
\usepackage[margin=0.86in]{geometry}
\usepackage{amsmath,amssymb,booktabs,array}
\usepackage[hidelinks]{hyperref}
\setlength{\parindent}{0pt}
\setlength{\parskip}{5pt}
\setlength{\emergencystretch}{2em}
\newcommand{\Cl}{\operatorname{Cl}}
\newcommand{\file}[1]{\texttt{\detokenize{#1}}}
\begin{document}
\begin{center}
{\Large\bfseries Imaginary quadratic fields of class-group exponent six}\\[5pt]
{\large A conditional classification with reproducible certificates}\\[5pt]
Research note for independent review \(\cdot\) 23 November 2026
\end{center}

\section*{1. The classification}
\textbf{Claim (assuming Dirichlet GRH).}
Exactly 432 imaginary quadratic fields have ideal class group of exponent six.
Their fundamental discriminants and class group structures are listed in
\file{known_exponent_six.csv}. The largest absolute discriminant is \(5\,761\,140\).

Let \(D<0\) be a fundamental discriminant, \(d=|D|\), and
\(E(D)=\exp\Cl(\mathcal O_{\mathbb Q(\sqrt D)})\).
Elsenhans, Kl\"uners and Nicolae (EKN) counted the 432 fields through
\(L=3.1\cdot10^{20}\). Their Theorem~3 and Table~2 give
\(Q_0=2.5\cdot10^{25}\) as an upper bound under the extended Riemann
hypothesis for quadratic fields. The relevant zeta functions factor as
\(\zeta(s)L(s,\chi_D)\). Aiken, Ross, Shvydkoy and Xue identify the same
interval \([L,Q_0]\) in Theorem~1.3(5). We use EKN's established low-range
classification and exclude every \(D\) in this interval with \(E(D)\mid6\).

\textbf{Norm equation.}
If \(p\mathcal O_K=\mathfrak p\overline{\mathfrak p}\) splits and \(E(D)\mid6\),
write \(\mathfrak p^6=(\alpha)\) with
\(\alpha=(x+y\sqrt D)/2\). A rational generator has equal valuations at the
two primes above \(p\), so \(y\ne0\). A generator of trace zero satisfies
\(\bar\alpha=-\alpha\), which would make its ideal equal to its conjugate.
After changing signs and conjugating if necessary,
\[
                 4p^6=x^2+dy^2,\qquad x,y\ge1 .
\]
For odd \(p\), reduction modulo \(p\) shows that \(p\mid x\) implies \(p\mid y\),
since \(p\nmid D\). Then \(\alpha\in p\mathcal O_K\) (use the integral basis
and \(p\) odd), contradicting the valuation at \(\bar{\mathfrak p}\).
In particular \(d<4p^6\). Comparing with \(L\) gives \(p\ge2069\) for every split prime.

\textbf{Finite exclusion.} The explicit formula on page 2 gives a split prime
below 3380. The least split prime therefore satisfies \(2069\le p\le3373\), and
\[
 d\le4\cdot3373^6=5\,890\,579\,727\,931\,596\,470\,756 .
\]
For all 165 possible primes our two implementations exhausted the necessary
trace conditions on page 3. Both found zero traces. This excludes the
entire interval and proves the conditional claim, subject to the supplied
finite computations. The arithmetic inputs, program sources and run records
are included for independent review.

\newpage
\section*{2. An explicit upper bound for the least split prime}
For \(q=d\) and \(X>2\), take \(g\) constant on each of 64 equal subintervals
of \([0,\log X)\), with the signed integer values in
\file{positive_kernel_verify.py}, and zero elsewhere. Its normalized
autocorrelation \(F(u)=\int g(t)g(t+u)\,dt/\int g^2\) is even, piecewise linear,
has \(F(0)=1\), and vanishes for \(|u|\ge\log X\). The Fourier transform
\(h(t)=\int F(u)e^{itu}\,du\) is nonnegative on the real axis. In every
fixed horizontal strip it decays as \(O_X((1+|\Re t|)^{-2})\).

For the odd primitive quadratic character \(\chi_D\), the Guinand--Weil
explicit formula (Carneiro et al., Lemma~9) now gives, under GRH,
\begin{align*}
2\sum_{n<X}\frac{\Lambda(n)\chi_D(n)}{\sqrt n}F(\log n)
 &\le \log(q/\pi)+\psi(3/4)+J_F,\\
J_F&=\int_0^\infty\frac{2e^{-3u/2}}{1-e^{-2u}}(1-F(u))\,du .
\end{align*}
Indeed the zero sum consists of nonnegative \(h(\gamma)\). Fourier inversion
in the integral representation of the digamma function gives \(J_F\).

If \(p<X\) is inert its contribution is
\(A_p=2\log p\,\sum_{p^k<X}(-1)^k F(k\log p)/p^{k/2}\); a ramified prime
contributes zero. If none splits, for every \(\lambda\ge0\) the left side is at least
\(\sum_{p<X}\min(A_p,\lambda\log p)-\lambda\log q\), since
\(\sum_{p\mid D}\log p\le\log q\). Using
\(\log\pi-\psi(3/4)>2.22\), the following positive margin forces a split prime:
\[
   \sum_{p<X}\min(A_p,\lambda\log p)-(1+\lambda)\log Q+2.22-J_{\rm up}>0,
   \qquad q\le Q,\quad J_{\rm up}\ge J_F .
\]
Every successful row permits \(Q\gets4X^6\) by the norm equation.
Starting at \(Q_0\), the outward-rounded lower bounds are
\begin{center}
\begin{tabular}{rrr@{\qquad\qquad}rrr}
\toprule \(X\)&\(\lambda\)&Margin&\(X\)&\(\lambda\)&Margin\\\midrule
4650&.0405&1.402389&3400&.0442&.181594\\
3750&.0428&1.155681&3390&.0443&.133120\\
3500&.0439&.444949&3385&.0443&.108959\\
3420&.0442&.207442&3380&.0443&.075457\\\bottomrule
\end{tabular}
\end{center}
Read down the left half, then the right. On the first cell of width
\(\Delta=(\log X)/64\), the weight of \(1-F\) is at most \(1/u\), giving
\(1-F(\Delta)\) as an upper bound for that integral. Each remaining cell
uses sixteen subintervals, multiplying the decreasing weight at the left
endpoint by the exact integral of \(1-F\). Beyond \(\log X\) one adds
\(4/[3X^{3/2}(1-X^{-2})]\). The code encloses the resulting finite
expression by intervals with 48-digit decimal arithmetic.

\textbf{Independent analytic certificate.} The logarithmic weight in
Lemmas~2.2 and 2.3 of Lamzouri, Li and Soundararajan also suffices.
Taking the residues in (2.8) gives
\(-E_1(y)\ge\gamma/2-(\gamma/2+\log2)/y\). At \(y=1020\) the six cutoffs
are \(4760,3895,3740,3708,3701,3700\).
Each successful row gives \(Q\gets4P^6\), with \(P\) the largest prime at
most the cutoff. This gives \(p_1\le3697\). Combining ramification costs
prime by prime and allowing a single exceptional split prime at \(X=3722\)
also gives \(p_2\le3719\). The corrected integer-interval certificate and
residue derivation are in \file{split_bound_proof.md} and its companion files.
The CPU run covers all 205 primes through 3697.

\newpage
\section*{3. Exhaustion of the trace interval}
For each candidate least split prime \(p\) define \(A=2p^3\) and
\(U_p=\lfloor\sqrt{A^2-L}\rfloor\). Any field in the gap supplies an integer
\(1\le x\le U_p\), \(p\nmid x\), and for every odd prime \(\ell<p\),
\[
 \left(\frac{x^2-A^2}{\ell}\right)
 =\left(\frac{D y^2}{\ell}\right)\in\{0,-1\}.
\]
If \(\ell\mid yD\), the symbol is zero and the residue is retained.
The condition at 2 and the fundamental-discriminant conditions can both be
omitted: this enlarges the set being exhausted. No integer survives even
these relaxed tests, so no factorization or class-group computation is
needed in the gap.

\begin{center}
\begin{tabular}{lrrr}
\toprule Computation&Prime buckets&Trace positions&Survivors\\ \midrule
CPU, through 3373&165&\(6\,140\,632\,911\,929\)&0\\
OpenCL, through 3373&165&\(6\,140\,632\,911\,929\)&0\\
CPU, through 3697&205&\(9\,649\,517\,432\,170\)&0\\ \bottomrule
\end{tabular}
\end{center}
Trace positions denote \(\sum_p U_p\) before congruence filtering.
The C\# sieve forms a CRT wheel for \(5,7,11,13,17,19\), intersects
64-bit periodic masks through 73 and checks the remaining primes on survivors.
The OpenCL program independently builds a different wheel through 13. It
filters on the device through 131 and sends the resulting 92,197 candidates
to a host-side exact residue test. The source, inclusive endpoints and
per-prime counts are in the bundle. The CPU self-test also compares masks
with Euler's criterion in intervals above \(7\cdot10^{10}\).

\textbf{Two-prime check.} The analytic certificates also bound a second
split prime. For instance, at \(X=3387\) a single exceptional split prime
costs at most \(0.148234\) in the explicit formula; the remaining margin is
greater than \(0.009267\). For two primes \(p<r\) with witnesses \((x,y),(u,v)\),
eliminating \(d\) gives
\[
(vx-yu)(vx+yu)=4(vp^3-yr^3)(vp^3+yr^3).
\]
Exhausting divisors for \(p\le3698,\ r\le3720\) gives \(6\,422\,835\)
primitive trace pairs before the odd-prime tests and zero after them.
\file{pair_method.md} derives the divisor range and covers the zero difference.

\textbf{Regression in the known range.} At \(L=1\) and \(p\le193\), two
implementations agreed on all 16,738 surviving traces across all 44 prime
buckets. Reconstruction of the fundamental discriminant and exact ideal-class
arithmetic recovered all 432 published fields with their class numbers.
For \(D=-5\,761\,140\) a witness is \(p=41,\ x=136498,\ y=8\), with class
number 576. The small-range run checks implementation behavior; the norm
equation and complete large-range run supply the exclusion.

\textbf{Sources.} EKN, \href{https://arxiv.org/abs/1803.02056}{arXiv:1803.02056},
Theorem~1, Lemma~1 and Table~2; \href{https://math.uni-paderborn.de/en/ag/ca/research/exponent/exponent0}{Paderborn class-group data}.
Aiken et al., \href{https://arxiv.org/abs/2605.26020}{arXiv:2605.26020},
Theorem~1.3(5). Lamzouri et al., \href{https://arxiv.org/abs/1309.3595}{arXiv:1309.3595},
Lemmas~2.2 and 2.3. Carneiro et al.,
\href{https://arxiv.org/abs/2404.08380}{arXiv:2404.08380}, Lemma~9.
\end{document}
